%%%%%%%%%%%%Outline%%%%%%%%%%%%%%%%%%

%%%
% message<RQ1: efficiency of generated native code>
% experiment 1.1: micro benchmark
% experiment 1.2: transform-family ablation study
% experiment 1.3: prove that native-code opportunities exist in real programs


%%%
% message<RQ2: application-level performance benefit>
% experiment 1: Runnable production-program subset
% experiment 2: End-to-end workload evaluation

%%%
% message<RQ3: overhead of online re-compilation>
% experiment 1: Re-compilation pipeline breakdown
% experiment 2: Steady-state daemon overhead
% experiment 3：Scaling and disruption under load
% experiment 4: Transparency

%%%
% message<RQ4: safety and correctness>
% experiment 1: Safety preservation
%  - how many programs are accepted / rejected by the verifier
%  - 
% experiment 2: Fail-safe rejection
% experiment 3: Semantic correctness

%%%
% message<RQ5: extensibility>
% experiment 1: Engineering cost of adding a new capability
% case study 1: New capability case study

% note
% benchmark setup, 完整的，see merlin 和 k2
% spec: detail 的 设计，具体做什么设计，要出什么图，代码需要什么修改 来 匹配 eval 
% 详细 描述 图是什么样子的，形式化的 定义
% 期待 的 图是什么样子

% note 
% safety boundary != optimization boundary


%%%%%%%%%%%%Outline%%%%%%%%%%%%%%%%%%

\section{Evaluation}
\label{sec:evaluation}

In this section, we aim to answer the following research questions:
\begin{itemize}
    \item \textbf{RQ1}: How effectively does \tool improve the efficiency of generated native code?
    \item \textbf{RQ2}: How much application-level performance benefit does \tool deliver in realistic deployments?
    \item \textbf{RQ3}: What is the overhead of online re-compilation in \tool?
    \item \textbf{RQ4}: Does \tool preserve the existing eBPF safety model and semantic correctness of its rewrites?
    \item \textbf{RQ5}: How easily can \tool be extended with new optimization capabilities?
\end{itemize}

We evaluate \tool along five complementary dimensions. 
%
RQ1 examines whether \tool improves the efficiency of the native code ultimately emitted by the kernel JIT, including both code compactness and low-level execution efficiency.
%
RQ2 asks whether these low-level code improvements translate into measurable application-level benefits in realistic deployments.
%
RQ3 measures the online cost of using \tool in practice, including the overheads of profiling, re-compilation, and program replacement on live systems.
%
RQ4 evaluates whether \tool preserves the existing eBPF safety model and whether its rewrites remain semantically consistent with the original programs.
%
RQ5 evaluates whether \tool is truly extensible by measuring the engineering effort needed to introduce new optimization capabilities.



\subsection{Evaluation Setup}
\label{subsec:evaluation-setup}

\noindent\textbf{Hardware.}


\noindent\textbf{Software.}


\noindent\textbf{Benchmarks.}
%
We use four benchmark layers. First, we use the active supported-family microbenchmarks to isolate improvements in generated native code. Second, we use a production-program corpus to measure both code-generation coverage and program-level performance on real BPF objects. Third, we use end-to-end deployments, with Tracee and Katran as the primary case studies for security/observability and networking, respectively. Fourth, we use negative and integration tests to evaluate fail-safe safety, semantic regression risk, and live-path stability. For RQ5, we additionally evaluate one newly added optimization capability as an extensibility case study.

\noindent\textbf{Evaluation Metrics.}
%
% Across the five RQs, we report native code size, native instruction count, execution time, cycles, instructions, and branch-miss counters for low-level measurements; throughput, latency, events per second, drop rate, and CPU utilization for application-level measurements; and end-to-end reJIT latency, background daemon cost, and transient disruption for online overhead. Unless a benchmark family requires a different convention, we report medians across repeated runs and show either 95\% confidence intervals or per-trial points for paired comparisons.

\subsection{RQ1: Efficiency of Generated Native Code}
\label{subsec:rq1}

RQ1 focuses on the quality of the native code ultimately emitted by the kernel JIT after \tool rewrites a live program. We therefore evaluate \tool first in a controlled microbenchmark setting, then summarize the improvement by transform family, and finally measure how often the same rewrite opportunities appear in real programs.

\noindent\textbf{Supported-Family Microbenchmarks.}
We first evaluate the subset of the active microbenchmark suite that corresponds to optimization families currently supported by the default \tool pipeline. Concretely, we include benchmarks that exercise conditional-select lowering, wide-memory-load simplification, and rotate-related idioms, and compare three execution modes: the stock kernel JIT, the same program after a live reJIT through \tool, and an optional stronger compiler-generated reference as an upper bound. For each benchmark, we record execution time, cycles, instructions, branch misses, and the size of the generated JIT image.

\revise{We expect the reJITed version to outperform the stock kernel JIT on most applicable microbenchmarks. The most stable code-size gains should come from wide-memory simplification and rotate-related rewrites, while conditional-select benchmarks should still improve on average but exhibit higher variance because their profitability depends more strongly on branch behavior and input distributions. If the stronger compiler reference is included, it will likely remain the best overall upper bound, but \tool should close a visible fraction of the remaining native-code gap.}

\noindent\textbf{Family-Sliced Summary.}
To ensure that the overall improvement is not driven by a single benchmark, we regroup the same microbenchmark results by transform family and report one aggregate slice per family. For each family, we summarize median speedup, median code-size reduction, and the main counter deltas relative to the stock kernel JIT.

\revise{We expect different families to contribute in different ways. Wide-memory rewrites should dominate code shrink, rotate-related rewrites should improve both compactness and execution efficiency, and conditional-select rewrites should show the strongest workload sensitivity. This breakdown should make clear that the benefit of \tool is not concentrated in one narrow code pattern.}

\noindent\textbf{Corpus Code-Generation Coverage and Code-Size Delta.}
Finally, we move from synthetic microbenchmarks to real programs. We run the production corpus in code-size mode, enumerate eligible rewrite sites, and measure the change in JIT image size after applying \tool. We report both the fraction of programs with at least one eligible site and the aggregate site distribution by subsystem and transform family.

\revise{We expect rewrite opportunities to appear across multiple subsystems rather than only in synthetic benchmarks. The eligible subset should show a positive median reduction in JIT image size, although the magnitude will likely be smaller and more heterogeneous than in the isolated microbenchmark suite. This experiment should establish that the native-code opportunities targeted by \tool are present in realistic programs.}

\subsection{RQ2: Application-Level Performance Benefit in Realistic Deployments}
\label{subsec:rq2}

RQ2 asks whether low-level native-code improvements translate into application-visible gains. We therefore evaluate \tool on real programs from the production corpus and on two end-to-end deployment scenarios: Tracee for security/observability and Katran for networking.

\noindent\textbf{Runnable Real-Program Subset from the Production Corpus.}
We first identify the subset of corpus programs that satisfy three conditions: they can be exercised reproducibly, both the stock and reJITed versions produce comparable measurements, and the program contains at least one eligible rewrite site. On this runnable subset, we measure per-program execution time and report the distribution of speedups achieved by \tool. Programs with no eligible sites, unstable triggers, or incomparable measurements are reported separately and excluded from the aggregate speedup statistics.

\revise{We expect the median speedup on the eligible subset to be positive, but materially smaller than on the synthetic microbenchmarks. This is the right trend: once helper costs, map accesses, and full application behavior are included, the impact of better generated native code should remain visible but become more modest. Reporting the excluded programs explicitly should also make the evaluation more credible by separating absence of opportunity from measurement instability.}

\noindent\textbf{Tracee End-to-End Deployment.}
We next evaluate \tool in a realistic security/observability deployment using Tracee. We use the current end-to-end harness to run representative workloads, including process creation, file I/O, and network activity, and compare the stock configuration against the same deployment after live reJIT. For each workload, we measure application throughput, Tracee event rate, drop counters, agent CPU utilization, and average BPF execution time.

\revise{We expect application-level gains to be smaller than the corresponding low-level BPF-time reductions but still directionally consistent. The most stable positive signal should be reduced average BPF time per event, while end-to-end throughput gains may vary by workload depending on how much of the total execution time is attributable to BPF processing. Drop counts should remain unchanged or near zero.}

\noindent\textbf{Katran End-to-End Deployment.}
We then evaluate \tool in a networking deployment using Katran. We drive a repeatable traffic workload, compare the stock path against the live reJITed path, and report request throughput, packet rate, tail latency, system CPU utilization, and average BPF execution time. To avoid order bias, the final evaluation should use a paired or counterbalanced design rather than a fixed stock-then-reJIT sequence.

\revise{We expect reduced average BPF time per packet to be the most consistent signal, with throughput or latency improvements that are positive but potentially modest. If the final setup removes phase-order bias and stabilizes the traffic generator, a small but repeatable application-level benefit should be observable.}

\subsection{RQ3: Overhead of Online Re-Compilation}
\label{subsec:rq3}

RQ3 measures the operational cost of applying \tool to live programs. We evaluate the online cost from three angles: the latency of invoking a reJIT request, the breakdown of time spent inside the reJIT pipeline, and the steady-state overhead of keeping the daemon enabled during normal operation.

\noindent\textbf{One-Shot Versus Persistent-Server ReJIT Requests.}
We first compare the two ways of invoking a live optimization. In one mode, the daemon is launched as a one-shot command for each optimization request. In the other, a persistent server process receives reJIT requests over its socket interface. We measure end-to-end latency for repeated optimizations of the same live target program and report medians and tail latency.

\revise{We expect persistent server mode to be clearly faster than repeated one-shot invocations, with most of the savings coming from avoiding process startup and repeated initialization. This comparison is important because it determines whether \tool should be treated as an occasional offline utility or as an always-on online service.}

\noindent\textbf{ReJIT Pipeline Breakdown.}
To understand where time is spent during a live reJIT, we instrument the pipeline and record timestamps around the main stages: selecting the target program, retrieving its original bytecode, analysis, rewrite, the verifier phase inside the reJIT path, the JIT phase, and completion of the atomic replacement. We run this experiment on both a small and a large target program.

\revise{We expect the verifier and JIT phases to dominate total latency, with analysis and rewriting accounting for a noticeably smaller fraction of the end-to-end cost. If this expectation holds, it would support the design choice of keeping optimization policy in userspace while relying on the kernel to enforce the existing verification and code-generation pipeline.}

\noindent\textbf{Steady-State Daemon Overhead.}
We next quantify the background cost of keeping \tool enabled during normal operation. On stable Tracee and Katran workloads, we compare four daemon states: daemon disabled, watch only, profile only, and the full adaptive path used in practice. We measure daemon CPU utilization, resident memory footprint, and any induced change in workload throughput or latency.

\revise{We expect watch-only and profile-only modes to impose near-negligible overhead, while the full online path should remain small relative to the application-level benefit recovered by reJIT. This result is necessary to justify \tool as a continuously running service rather than a manually triggered optimization pass.}

\noindent\textbf{Disruption During a Live ReJIT Event.}
Finally, we study whether one live reJIT causes visible disturbance to an ongoing workload. We run a steady Tracee or Katran workload, trigger exactly one reJIT event in the middle of the run, and collect a short time series centered on the event. We then report the reJIT completion time, the maximum transient throughput dip, the maximum latency spike, and any transient packet or event loss.

\revise{We expect any disturbance to be short and small, with no catastrophic interruption of the workload. A brief and bounded transient would be consistent with the design goal of supporting live optimization without tearing down the running application stack.}

\subsection{RQ4: Safety Preservation and Semantic Correctness}
\label{subsec:rq4}

RQ4 separates two questions that are often conflated. Safety is a strong systems guarantee: rewritten programs must preserve the existing eBPF safety model because they are re-verified through the same kernel mechanism as ordinary loads. Semantic correctness is instead an empirical property of the user-space rewrite logic: we must provide evidence that the rewritten program behaves like the original program on the inputs and states that matter.

\noindent\textbf{Targeted Adversarial Fail-Safe Safety Test.}
We first run the targeted adversarial reJIT test suite, which crafts invalid rewrite payloads that attempt to violate memory safety, control-flow integrity, type safety, privilege separation, and other verifier-enforced invariants. For each category, we record whether the reJIT path rejects the payload, whether the original live program remains runnable afterward, and whether the kernel reports any warning, bug, or crash.

\revise{We expect all invalid adversarial rewrites to be rejected, the original program to remain unchanged, and the kernel to report no warnings, oopses, or panics. This is the clearest test of the claim that a buggy or compromised userspace optimizer can at worst fail closed, rather than compromising the kernel's safety model.}

\noindent\textbf{Random Fuzzed-Bytecode Fail-Safe Test.}
We then stress the same fail-safe claim with randomized and mutated bytecode. We run the existing fuzzing modes for a large number of rounds, record acceptance and rejection outcomes, and verify that failed attempts do not corrupt or replace the original live program. Any accepted random program is manually inspected to confirm that it remains verifier-safe.

\revise{We expect the overwhelming majority of malformed programs to be rejected and, critically, to leave the original program unaffected. Even if a small fraction of mutated inputs are accepted, the important property is that the verifier continues to enforce safety and that \tool does not bypass the kernel's existing protection boundary.}

\noindent\textbf{Live Integration Safety Smoke Test.}
Beyond negative tests, we also validate the normal live path. We run the integration test that enumerates a live program, rewrites it, applies the rewritten version, and then exercises the bulk-apply path, while checking the kernel log for warnings. This experiment confirms that the intended ``happy path'' remains stable on a real loaded program.

\revise{We expect all phases of the live path to succeed without kernel warnings. Together with the adversarial and fuzz tests, this experiment should show that \tool is safe both when it is used correctly and when it is driven with malformed inputs.}

\noindent\textbf{Differential Semantic Testing.}
To evaluate semantic correctness, we compare the original and rewritten versions of the same program on identical inputs and identical initial state. For each supported rewrite family, we compare return values, map updates, emitted events, and other externally visible side effects. We use the microbenchmarks for family-specific coverage and at least one real corpus program for each supported family whenever possible.

\revise{We expect no semantic mismatches on the supported rewrite families. More importantly, this experiment should make the boundary of the claim explicit: the verifier provides safety, while semantic preservation is supported by systematic differential testing rather than asserted as a theorem.}

\subsection{RQ5: Extensibility of the Optimization Framework}
\label{subsec:rq5}

RQ5 evaluates whether \tool is genuinely extensible rather than merely a fixed bundle of hard-coded rewrites. We therefore add one new optimization capability to the current system, measure the implementation effort required to do so, and then demonstrate that the new capability delivers real value while preserving the same operational model.

\noindent\textbf{Engineering Cost of Adding One New Optimization Capability.}
We select one rewrite family that is already tracked by the benchmark suite but is not yet part of the default optimization pipeline, and we implement it as a new \tool capability. We then measure the engineering effort required to add the feature, including lines of code and touched files in the userspace daemon, any optional architecture-specific module support, and any changes to the core kernel or verifier.

\revise{We expect the vast majority of the implementation effort to remain in the daemon and, if needed, in the optional architecture-specific module, with no changes or only negligible changes to the core kernel and verifier. This would directly support the claim that \tool externalizes optimization policy without repeatedly perturbing the kernel's trusted core.}

\noindent\textbf{Benefit of the New Capability.}
To show that the new capability is not only easy to add but also useful, we evaluate it on both a dedicated microbenchmark and at least one real corpus program that exposes the same opportunity. We report the number of matched sites, the resulting code-size change, and the corresponding runtime improvement.

\revise{We expect the dedicated microbenchmark to show a clear win and the real program to show at least one measurable positive effect, even if the gain is smaller. This experiment is important because it demonstrates that extensibility in \tool is not merely an implementation convenience; new capabilities can be added and immediately exercised through the same live reJIT workflow.}

\noindent\textbf{Module Lifecycle and Graceful Degradation.}
Finally, we test the lifecycle behavior of the new capability by comparing three states: the relevant module or capability is available, it is absent, and it is removed after deployment. We verify that the system either safely declines to apply the optimization or degrades gracefully, without breaking live programs or weakening the safety guarantees established in RQ4.

\revise{We expect capability absence or module unload to remove the optimization opportunity rather than to break the workload. A safe fallback or safe rejection path would show that extensibility in \tool remains robust even outside the ideal happy path.}